\chapter{Conclusiones}

En el presente trabajo de memoria se cumplió con el objetivo general de implementar y evaluar algoritmos de detección de objetos, basados en inteligencia artificial, para la identificación automática de fallas en correas transportadoras. Asimismo, los objetivos específicos y alcances propuestos fueron logrados con éxito. Se implementaron y evaluaron dos arquitecturas convolucionales asociadas a modelos de detección de objetos de una sola etapa (YOLOv5 y SSD) junto con modelos basados en mecanismos de atención global con arquitectura (\textit{transformer}) (DETR y DINO). Todo el proceso de entrenamiento, validación y prueba se realizó utilizando una base de datos pública debidamente estructurada, considerando cinco clases de falla y un esquema estandarizado que permitió una evaluación cuantitativa ordenada y precisa.\\


Respecto a los resultados y discusiones, la variante \textit{small} (s) de YOLOv5 se consolidó como el modelo superior al maximizar la eficiencia, mientras que SSD512 obtuvo el rendimiento más deficiente. El mayor desafío recayó en la detección deficiente de la clase perforación (\textit{Puncture}); no obstante, DINO-R50-4scale demostró ligera superioridad prediciendo este defecto gracias a su atención global y su mecanismo mejorado de eliminación de ruido contrastivo. \\

El desempeño de las arquitecturas \textit{Transformers} se vio penalizado por la escala de los datos, la brecha de dominio y las restricciones físicas. Su extremo consumo de memoria forzó a reducir los tamaño de lote, ralentizando la convergencia de los gradientes y aumentando los tiempos de entrenamiento. Al mantener un régimen estandarizado sin optimización profunda de hiperparámetros, los detectores convolucionales \textit{One-Stage} se ratifican como la alternativa de mayor viabilidad práctica e industrial en la actualidad. \\

Como trabajo futuro se propone: una optimización de hiperparámetros para cada modelo; la ejecución de entrenamientos en escritorios fijos dedicado, en vez de equipos portátiles; y la implementación de arquitecturas híbridas (como RT-DETR o DELO), orientadas a fusionar la eficiencia de extracción local de las CNN con el mecanismo de atención global de los \textit{Transformers}.